\chapter*{Conventions}
\addcontentsline{toc}{chapter}{Conventions}
\markboth{Conventions}{Conventions}

These conventions are global.  A chapter that departs from one of them says
so before the first departure.

\section*{Sets, topology, and extended real numbers}

The set of natural numbers is \(\mathbf N=\{0,1,2,\ldots\}\).  The extended
real line is \(\Rbar=\mathbf R\cup\{-\infty,+\infty\}\).  We use the order
topology when both endpoints occur and the evident subspace topology on
one-sided extensions.

A compact space is not assumed metrizable.  Convergence statements are
therefore formulated for nets unless a sequential reduction has been proved.
A subnet is a Willard subnet: its indexing map is order preserving and
cofinal.  If \((f_i)\) is a family of extended-real functions, then
\((\sup_i f_i)^*\) denotes the least upper semicontinuous majorant of its
pointwise supremum.

We never form \((+\infty)-(+\infty)\) or
\((-\infty)-(-\infty)\).  A displayed identity involving extended-real
functions is understood only on the locus on which both sides are defined,
unless it is formulated monotonically without subtraction.  The convention
\(0\cdot(+\infty)=0\) is used only inside integrals after it has been justified
by monotone truncation.

All measures are positive Radon measures unless signed measures are
explicitly mentioned.  Weak convergence means convergence against continuous
real-valued test functions on the ambient compact Hausdorff space.

\section*{Fields, absolute values, and valuations}

A non-Archimedean field \(K\) is equipped with a multiplicative absolute value
\(|\cdot|:K\to\mathbf R_{\ge0}\) satisfying the ultrametric inequality.  It is
complete unless otherwise stated.  The valuation associated to the absolute
value is
\[
  v_K(a)=-\log|a|.
\]
Thus larger valuation means smaller norm.  When the absolute value is trivial,
\(v_K(a)=0\) for every \(a\in K^\times\).

A semivaluation on a ring \(A\) is a map
\(v:A\to[0,+\infty]\) satisfying
\[
 \begin{gathered}
 v(0)=+\infty,\qquad v(1)=0,\\
 v(fg)=v(f)+v(g),\qquad
 v(f+g)\ge\min\{v(f),v(g)\}.
 \end{gathered}
\]
Its support is the prime ideal
\(\{f\in A\mid v(f)=+\infty\}\).  For an additive valuation with values in a
general ordered group, nonnegativity is not imposed unless a center has been
specified.

\section*{Algebraic and analytic spaces}

A variety over a field is a separated integral scheme of finite type, except
in chapters that explicitly allow reduced or reducible schemes.  Projective
means projective over the ground field.  A model of a proper \(K\)-variety is
specified in the chapter at hand: algebraic models over the valuation ring
and admissible formal models are not silently identified.

For a \(K\)-scheme \(X\), its Berkovich analytification is denoted
\(X^{\An}\).  If \(K\) is trivially valued, points of \(X^{\An}\) include
semivaluations with nonzero support; they are not discarded.  A divisorial
point is normalized relative to the ground valuation or the scaling action
specified in the relevant chapter.

A complex space is reduced, paracompact, and Hausdorff.  A compact Kähler
space is assumed irreducible when a function field or a generic Lelong number
is used.  Modifications are proper bimeromorphic morphisms.

\section*{Metrics and potentials}

On a non-Archimedean analytic line bundle we use additive metric weights.  If
\(\|\cdot\|_\phi\) and \(\|\cdot\|_{\phi_0}\) are metrics on the same line
bundle, their relative weight is the function \(\phi-\phi_0\) determined by
\[
  \|s(x)\|_\phi=\rme^{-(\phi-\phi_0)(x)}\|s(x)\|_{\phi_0}.
\]
Changing the reference metric therefore adds a continuous function.  The
notations model metric, formal metric, PL metric, and Fubini--Study metric are
defined separately in \Cref{chap:metrics}.

On a complex line bundle we follow the convention of the supplied book: if
\(h\) and \(h_0\) are Hermitian metrics, their relative psh potential
\(\varphi\) is determined by
\[
  |s|_h^2=\rme^{-\varphi}|s|_{h_0}^2 .
\]
Thus the additive weight of the underlying norm is \(\varphi/2\).  This
factor of two is displayed whenever complex and non-Archimedean metric
weights are compared.

The sign of model functions is fixed as follows.  If \(D\) is a vertical
Cartier divisor, locally cut out by \(f=0\), then \(\cO(D)\) is locally
generated by \(f^{-1}\), and
\[
  \varphi_D=\log|\cO(D)|=\log|f^{-1}|=-\log|f|=v_x(f).
\]
Thus effective vertical divisors give nonnegative model functions.  In the
discretely valued normalization \(v_K(\pi)=1\), one has
\(\varphi_{\cX_s}=1\).  These two tests determine every later sign involving
vertical divisors.

The symbol \(\PSH(X,\theta)\) always includes the ambient space and the class
or form in its definition.  It may refer to CLD, BFJ, BJ, or complex psh
objects only after the chapter has declared which category is in force.  We
never infer equality of two psh classes from their notation alone.

\section*{Differential operators and curvature}

On a complex space we use
\[
  \ddc=\frac{\rmi}{2\pi}\partial\bar\partial.
\]
Thus the Poincaré--Lelong formula for a nonzero meromorphic function is
\[
 \ddc\log|f|^2=[\operatorname{div}(f)],
 \qquad
 \ddc\log|f|=\tfrac12[\operatorname{div}(f)],
\]
where the second equality is the equivalent half-logarithm convention.

In the Chambert--Loir--Ducros calculus the operators are denoted
\(\dprime\), \(\ddprime\), and \(\dprimeddprime\).  Their sign and lattice
normalizations are not forced into the complex notation.  \Cref{app:normalizations} records
the conversion used in each comparison theorem.

\section*{Monge--Ampère measures and energy}

Let \(X\) be proper of pure dimension \(n\), and let \(L\) be an ample line
bundle.  Unless the word \emph{probability} appears, the Monge--Ampère measure
of an \(L\)-psh metric has total mass
\[
  (L^n).
\]
For a numerical class \(\theta\), the corresponding total mass is written
\(\vol(\theta)\) in the big case.  Probability normalization divides by this
number and is never used when the number may vanish.

Energy functionals are normalized to vanish at the chosen reference metric.
Changing that reference adds the appropriate affine constant; all variational
statements are invariant under this change.  The finite-energy class is
denoted \(\cE^1\), with an ambient subscript or superscript when more than one
theory is present.

\section*{Positivity words}

Positive means nonnegative unless \emph{strictly positive} is written.
Semipositive, nef, ample, big, Kähler, and pseudoeffective are not synonyms.
For a model line bundle, semipositivity is defined by relative nefness.  For a
continuous metric it is defined by an approximation theorem.  For a numerical
class it is defined in the relevant Néron--Severi or Bott--Chern vector space.

\section*{Citations and provenance}

When a theorem is quoted from the literature, its first reliable source is
named in the theorem heading or immediately before the statement.  A later
extension is stated separately.  Preprints are identified by version and
date.  Results from the supplied envelope manuscript are proved in full in
\cref{part:bj} and are not attributed to the earlier BJ addendum.

The Stacks Project is cited by tag.  \Cref{app:stacks} lists the exact statement used
at each tag, the hypotheses needed here, and the first place of application.
